\chapter{Bloch Geometry, Berry Holonomy, and the Spinor Sign}
\label{ch:bloch}

\section{Two-level state}
The square-root embedding becomes a normalized complex two-level state after introducing a relative phase:
\begin{equation}
\ket{\psi(\sigma,\phi)}
=\sqrt{\sigma}\ket{0}
+e^{i\phi}\sqrt{1-\sigma}\ket{1}.
\label{eq:qubit-state}
\end{equation}
Using \eqref{eq:sigma-theta}, this can be written
\begin{equation}
\ket{\psi(\theta,\phi)}
=\cos\!\left(\frac\theta2\right)\ket{0}
+e^{i\phi}\sin\!\left(\frac\theta2\right)\ket{1}.
\end{equation}
Up to an irrelevant global phase, this is the standard pure-qubit parametrization. Its Bloch vector is
\begin{equation}
\boldsymbol{n}=(\sin\theta\cos\phi,\sin\theta\sin\phi,\cos\theta).
\end{equation}
Since $\cos\theta=2\sigma-1$, the complement $\sigma\mapsto1-\sigma$ reflects the Bloch vector through the equatorial plane.

\begin{proposition}[Balanced point equals the Bloch equator]
The information-balanced condition $\sigma=1/2$ is equivalent to
\begin{equation}
n_z=0,
\end{equation}
so the full phase orbit lies on the Bloch equator.
\end{proposition}

\begin{proof}
From $n_z=\cos\theta=2\sigma-1$, one has $n_z=0$ if and only if $\sigma=1/2$.
\end{proof}

\section{Berry connection}
Let $\phi$ vary adiabatically around a full cycle while $\sigma$ is held fixed. Differentiate \eqref{eq:qubit-state}:
\begin{equation}
\partial_\phi\ket{\psi}
=i e^{i\phi}\sqrt{1-\sigma}\ket{1}.
\end{equation}
The Berry connection along the cycle is
\begin{equation}
\mathcal A_\phi
=i\braket{\psi}{\partial_\phi\psi}
=i\left[i(1-\sigma)\right]
=-(1-\sigma).
\end{equation}
Therefore the geometric phase is
\begin{equation}
\gamma_B(\sigma)
=\int_0^{2\pi}\mathcal A_\phi\,\dd\phi
=-2\pi(1-\sigma).
\label{eq:berry-sigma}
\end{equation}
A different gauge may shift $\gamma_B$ by an integer multiple of $2\pi$, but the holonomy
\begin{equation}
\mathcal U_B(\sigma)=e^{i\gamma_B(\sigma)}
\end{equation}
is gauge invariant.

The same result follows from the solid-angle formula. The latitude at polar angle $\theta$ encloses solid angle
\begin{equation}
\Omega=2\pi(1-\cos\theta),
\end{equation}
and a spin-$1/2$ state acquires
\begin{equation}
\gamma_B=-\frac{\Omega}{2}
=-\pi(1-\cos\theta)
=-2\pi(1-\sigma).
\end{equation}

\section{Equatorial holonomy}
\begin{theorem}[Information-balanced spinor holonomy]
For the state family \eqref{eq:qubit-state}, a full azimuthal cycle at the balanced point has
\begin{equation}
\gamma_B(1/2)=-\pi\pmod{2\pi},
\qquad
\mathcal U_B(1/2)=-1.
\end{equation}
\end{theorem}

\begin{proof}
Substituting $\sigma=1/2$ into \eqref{eq:berry-sigma} gives $\gamma_B=-\pi$. Euler's formula then gives $e^{-i\pi}=-1$.
\end{proof}

Thus the same point that maximizes binary entropy also produces the nontrivial sign holonomy of a spinor-like two-state cycle.

\section{Two pi and four pi}
For a spin-$1/2$ representation, a physical rotation by angle $\alpha$ about a unit axis $\boldsymbol{n}$ is represented by
\begin{equation}
U(\alpha)=\exp\!\left(-\frac{i\alpha}{2}\boldsymbol{n}\cdot\boldsymbol{\sigma}\right),
\end{equation}
where $\boldsymbol{\sigma}$ are the Pauli matrices. Therefore
\begin{equation}
U(2\pi)=-I,
\qquad
U(4\pi)=I.
\end{equation}
The factor $1/2$ is the representation weight responsible for the double cover $SU(2)\to SO(3)$. In the present construction, an analogous half-weight appears as the balance coordinate of the two complementary amplitudes.

It is essential to distinguish two claims:
\begin{enumerate}
\item \emph{Established:} spin-$1/2$ representations have a $-1$ sign under a $2\pi$ rotation.
\item \emph{Proposed identification:} the real coordinate of a zeta zero is represented by the population weight of a two-level state whose equatorial holonomy implements the zeta phase balance.
\end{enumerate}
The second is not implied merely by the first.

\section{Euler phase versus geometric phase}
The equation $e^{i\pi}=-1$ is algebraically universal. It does not by itself select $1/2$. The selection occurs because the Berry phase depends on the population coordinate:
\begin{equation}
\gamma_B(\sigma)=-2\pi(1-\sigma).
\end{equation}
Requiring the holonomy to equal $-1$ gives
\begin{equation}
-2\pi(1-\sigma)=\pi\pmod{2\pi},
\end{equation}
which yields
\begin{equation}
\sigma=\frac12\pmod1.
\end{equation}
In the probability interval $0\leq\sigma\leq1$, the nontrivial interior solution is $1/2$. The phase requirement therefore selects the same point as Shannon balance.

\begin{figure}[ht]
\centering
\includegraphics[width=0.66\textwidth]{bloch_equator.pdf}
\caption{Schematic Bloch geometry. Complement exchange reflects north and south populations, while $\sigma=1/2$ is the fixed equator.}
\label{fig:bloch}
\end{figure}
